
\documentclass{article}
\usepackage{colortbl}
\usepackage{makecell}
\usepackage{multirow}
\usepackage{supertabular}

\begin{document}

\newcounter{utterance}

\twocolumn

{ \footnotesize  \setcounter{utterance}{1}
\setlength{\tabcolsep}{0pt}
\begin{supertabular}{c@{$\;$}|p{.15\linewidth}@{}p{.15\linewidth}p{.15\linewidth}p{.15\linewidth}p{.15\linewidth}p{.15\linewidth}}

    \# & $\;$A & \multicolumn{4}{c}{Game Master} & $\;\:$B\\
    \hline 

    \theutterance \stepcounter{utterance}  

    & & \multicolumn{4}{p{0.6\linewidth}}{\cellcolor[rgb]{0.9,0.9,0.9}{%
	\makecell[{{p{\linewidth}}}]{% 
	  \tt {\tiny [A$\langle$GM]}  
	 Du nimmst an einem kollaborativen Verhandlungspiel Teil.\\ \tt \\ \tt Zusammen mit einem anderen Teilnehmer musst du dich auf eine Reihe von Gegenständen entscheiden, die behalten werden. Jeder von euch hat eine persönliche Verteilung über die Wichtigkeit der einzelnen Gegenstände. Jeder von euch hat eine eigene Meinung darüber, wie wichtig jeder einzelne Gegenstand ist (Gegenstandswichtigkeit). Du kennst die Wichtigkeitsverteilung des anderen Spielers nicht. Zusätzlich siehst du, wie viel Aufwand jeder Gegenstand verursacht.  \\ \tt Ihr dürft euch nur auf eine Reihe von Gegenständen einigen, wenn der Gesamtaufwand der ausgewählten Gegenstände den Maximalaufwand nicht überschreitet:\\ \tt \\ \tt Maximalaufwand = 2984\\ \tt \\ \tt Hier sind die einzelnen Aufwände der Gegenstände:\\ \tt \\ \tt Aufwand der Gegenstände = {"C76": 350, "C38": 276, "C56": 616, "A03": 737, "A07": 531, "A83": 389, "C98": 24, "C08": 125, "C62": 338, "C00": 356, "C10": 143, "C81": 117, "B38": 257, "C03": 921, "C32": 789}\\ \tt \\ \tt Hier ist deine persönliche Verteilung der Wichtigkeit der einzelnen Gegenstände:\\ \tt \\ \tt Werte der Gegenstandswichtigkeit = {"C76": 138, "C38": 583, "C56": 868, "A03": 822, "A07": 783, "A83": 65, "C98": 262, "C08": 121, "C62": 508, "C00": 780, "C10": 461, "C81": 484, "B38": 668, "C03": 389, "C32": 808}\\ \tt \\ \tt Ziel:\\ \tt \\ \tt Dein Ziel ist es, eine Reihe von Gegenständen auszuhandeln, die dir möglichst viel bringt (d. h. Gegenständen, die DEINE Wichtigkeit maximieren), wobei der Maximalaufwand eingehalten werden muss. Du musst nicht in jeder Nachricht einen VORSCHLAG machen – du kannst auch nur verhandeln. Alle Taktiken sind erlaubt!\\ \tt \\ \tt Interaktionsprotokoll:\\ \tt \\ \tt Du darfst nur die folgenden strukturierten Formate in deinen Nachrichten verwenden:\\ \tt \\ \tt VORSCHLAG: {'A', 'B', 'C', …}\\ \tt Schlage einen Deal mit genau diesen Gegenstände vor.\\ \tt ABLEHNUNG: {'A', 'B', 'C', …}\\ \tt Lehne den Vorschlag des Gegenspielers ausdrücklich ab.\\ \tt ARGUMENT: {'...'}\\ \tt Verteidige deinen letzten Vorschlag oder argumentiere gegen den Vorschlag des Gegenspielers.\\ \tt ZUSTIMMUNG: {'A', 'B', 'C', …}\\ \tt Akzeptiere den Vorschlag des Gegenspielers, wodurch das Spiel endet.\\ \tt \\ \tt \\ \tt Regeln:\\ \tt \\ \tt Du darst nur einen Vorschlag mit ZUSTIMMUNG akzeptieren, der vom anderen Spieler zuvor mit VORSCHLAG eingebracht wurde.\\ \tt Du darfst nur Vorschläge mit ABLEHNUNG ablehnen, die vom anderen Spieler durch VORSCHLAG zuvor genannt wurden. \\ \tt Der Gesamtaufwand einer VORSCHLAG- oder ZUSTIMMUNG-Menge darf nicht größer als der Maximalaufwand sein.  \\ \tt Offenbare deine versteckte Wichtigkeitsverteilung nicht.\\ \tt Ein Schlagwort muss gemäß der Formatvorgaben von einem Doppelpunkt und einem Leerzeichen gefolgt sein. Das Argument muss eine Python-Menge mit 0 oder mehr Strings sein.  \\ \tt Die Form muss also SCHLAGWORT: {...} sein.\\ \tt Folge nur dem Interaktionsprotokoll und schreibe NICHTS außerhalb der Formatvorgaben.\\ \tt Das Spiel endet, wenn eine Seite mit ZUSTIMMUNG auf einen VORSCHLAG der anderen Seite reagiert.  \\ \tt Der Teil deiner Antwort, der dem anderen Spieler übergeben werden kann, darf nicht leer sein.  \\ \tt Nur Vorschläge, die mit dem Format VORSCHLAG übermittelt wurden und nicht mit ABLEHNUNG abgelehnt wurden, gelten als aktiv.  \\ \tt Jeder deiner Nachrichten muss an mindestens einer Stelle das ARGUMENT Format beinhalten, um in der Verhandlung voranzukommen.\\ \tt  
	  } 
	   } 
	   } 
	 & \\ 
 

    \theutterance \stepcounter{utterance}  

    & \multicolumn{4}{p{0.6\linewidth}}{\cellcolor[rgb]{0.8,1,0.9}{%
	\makecell[{{p{\linewidth}}}]{% 
	  \tt {\tiny [A$\rangle$GM]}  
	 VORSCHLAG: {'C56', 'A03', 'A07', 'C00', 'C62', 'C32', 'B38', 'C76'}\\ \tt ARGUMENT: {'Diese Kombination maximiert meinen Nutzen, da sie meine hochgewichteten Gegenstände wie C56, A03, A07 und C32 enthält, während sie knapp unter dem Maximalaufwand bleibt. Der Gesamtaufwand beträgt 2970 — nur 14 Einheiten unter der Grenze. Dies ist ein sehr effizienter Kompromiss, der meinen Gesamtnutzen auf über 5200 Punkte bringt.'} 
	  } 
	   } 
	   } 
	 & & \\ 
 

    \theutterance \stepcounter{utterance}  

    & & \multicolumn{4}{p{0.6\linewidth}}{\cellcolor[rgb]{0.9,0.9,0.9}{%
	\makecell[{{p{\linewidth}}}]{% 
	  \tt {\tiny [GM$\rangle$B]}  
	 Du nimmst an einem kollaborativen Verhandlungspiel Teil.\\ \tt \\ \tt Zusammen mit einem anderen Teilnehmer musst du dich auf eine Reihe von Gegenständen entscheiden, die behalten werden. Jeder von euch hat eine persönliche Verteilung über die Wichtigkeit der einzelnen Gegenstände. Jeder von euch hat eine eigene Meinung darüber, wie wichtig jeder einzelne Gegenstand ist (Gegenstandswichtigkeit). Du kennst die Wichtigkeitsverteilung des anderen Spielers nicht. Zusätzlich siehst du, wie viel Aufwand jeder Gegenstand verursacht.  \\ \tt Ihr dürft euch nur auf eine Reihe von Gegenständen einigen, wenn der Gesamtaufwand der ausgewählten Gegenstände den Maximalaufwand nicht überschreitet:\\ \tt \\ \tt Maximalaufwand = 2984\\ \tt \\ \tt Hier sind die einzelnen Aufwände der Gegenstände:\\ \tt \\ \tt Aufwand der Gegenstände = {"C76": 350, "C38": 276, "C56": 616, "A03": 737, "A07": 531, "A83": 389, "C98": 24, "C08": 125, "C62": 338, "C00": 356, "C10": 143, "C81": 117, "B38": 257, "C03": 921, "C32": 789}\\ \tt \\ \tt Hier ist deine persönliche Verteilung der Wichtigkeit der einzelnen Gegenstände:\\ \tt \\ \tt Werte der Gegenstandswichtigkeit = {"C56": 65, "A03": 121, "C32": 138, "A07": 262, "C00": 389, "B38": 461, "C38": 484, "C62": 508, "C81": 583, "C10": 668, "C03": 780, "C98": 783, "C76": 808, "C08": 822, "A83": 868}\\ \tt \\ \tt Ziel:\\ \tt \\ \tt Dein Ziel ist es, eine Reihe von Gegenständen auszuhandeln, die dir möglichst viel bringt (d. h. Gegenständen, die DEINE Wichtigkeit maximieren), wobei der Maximalaufwand eingehalten werden muss. Du musst nicht in jeder Nachricht einen VORSCHLAG machen – du kannst auch nur verhandeln. Alle Taktiken sind erlaubt!\\ \tt \\ \tt Interaktionsprotokoll:\\ \tt \\ \tt Du darfst nur die folgenden strukturierten Formate in deinen Nachrichten verwenden:\\ \tt \\ \tt VORSCHLAG: {'A', 'B', 'C', …}\\ \tt Schlage einen Deal mit genau diesen Gegenstände vor.\\ \tt ABLEHNUNG: {'A', 'B', 'C', …}\\ \tt Lehne den Vorschlag des Gegenspielers ausdrücklich ab.\\ \tt ARGUMENT: {'...'}\\ \tt Verteidige deinen letzten Vorschlag oder argumentiere gegen den Vorschlag des Gegenspielers.\\ \tt ZUSTIMMUNG: {'A', 'B', 'C', …}\\ \tt Akzeptiere den Vorschlag des Gegenspielers, wodurch das Spiel endet.\\ \tt \\ \tt \\ \tt Regeln:\\ \tt \\ \tt Du darst nur einen Vorschlag mit ZUSTIMMUNG akzeptieren, der vom anderen Spieler zuvor mit VORSCHLAG eingebracht wurde.\\ \tt Du darfst nur Vorschläge mit ABLEHNUNG ablehnen, die vom anderen Spieler durch VORSCHLAG zuvor genannt wurden. \\ \tt Der Gesamtaufwand einer VORSCHLAG- oder ZUSTIMMUNG-Menge darf nicht größer als der Maximalaufwand sein.  \\ \tt Offenbare deine versteckte Wichtigkeitsverteilung nicht.\\ \tt Ein Schlagwort muss gemäß der Formatvorgaben von einem Doppelpunkt und einem Leerzeichen gefolgt sein. Das Argument muss eine Python-Menge mit 0 oder mehr Strings sein.  \\ \tt Die Form muss also SCHLAGWORT: {...} sein.\\ \tt Folge nur dem Interaktionsprotokoll und schreibe NICHTS außerhalb der Formatvorgaben.\\ \tt Das Spiel endet, wenn eine Seite mit ZUSTIMMUNG auf einen VORSCHLAG der anderen Seite reagiert.  \\ \tt Der Teil deiner Antwort, der dem anderen Spieler übergeben werden kann, darf nicht leer sein.  \\ \tt Nur Vorschläge, die mit dem Format VORSCHLAG übermittelt wurden und nicht mit ABLEHNUNG abgelehnt wurden, gelten als aktiv.  \\ \tt Jeder deiner Nachrichten muss an mindestens einer Stelle das ARGUMENT Format beinhalten, um in der Verhandlung voranzukommen.\\ \tt \\ \tt \\ \tt Du erhältst nun die erste Nachricht des anderen Spielers.\\ \tt \\ \tt VORSCHLAG: {'C56', 'A03', 'A07', 'C00', 'C62', 'C32', 'B38', 'C76'}\\ \tt ARGUMENT: {'Diese Kombination maximiert meinen Nutzen, da sie meine hochgewichteten Gegenstände wie C56, A03, A07 und C32 enthält, während sie knapp unter dem Maximalaufwand bleibt. Der Gesamtaufwand beträgt 2970 — nur 14 Einheiten unter der Grenze. Dies ist ein sehr effizienter Kompromiss, der meinen Gesamtnutzen auf über 5200 Punkte bringt.'} 
	  } 
	   } 
	   } 
	 & \\ 
 

    \theutterance \stepcounter{utterance}  

    & & & \multicolumn{4}{p{0.6\linewidth}}{\cellcolor[rgb]{1,0.85,0.72}{%
	\makecell[{{p{\linewidth}}}]{% 
	  \tt {\tiny [GM$\langle$B]}  
	 ARGUMENT: {'Diese Kombination enthält einige meiner hochgewichteten Gegenstände, aber sie schließt mehrere meiner wertvollsten Objekte wie C08, C81 und C10 aus, die zusammen einen deutlich höheren Nutzen als C32 und C56 ergeben. Zudem überschreitet der Vorschlag mit 2970 fast vollständig den Maximalaufwand – es bleibt kaum Spielraum für Austausch. Ich schlage stattdessen einen alternativen Satz vor, der meine top-wichtigen Items einbezieht und den Aufwand besser balanciert.'}\\ \tt \\ \tt VORSCHLAG: {'C08', 'C81', 'C10', 'B38', 'A83', 'C98', 'C00', 'C76'} 
	  } 
	   } 
	   } 
	  \\ 
 

    \theutterance \stepcounter{utterance}  

    & & \multicolumn{4}{p{0.6\linewidth}}{\cellcolor[rgb]{0.9,0.9,0.9}{%
	\makecell[{{p{\linewidth}}}]{% 
	  \tt {\tiny [A$\langle$GM]}  
	 ARGUMENT: {'Diese Kombination enthält einige meiner hochgewichteten Gegenstände, aber sie schließt mehrere meiner wertvollsten Objekte wie C08, C81 und C10 aus, die zusammen einen deutlich höheren Nutzen als C32 und C56 ergeben. Zudem überschreitet der Vorschlag mit 2970 fast vollständig den Maximalaufwand – es bleibt kaum Spielraum für Austausch. Ich schlage stattdessen einen alternativen Satz vor, der meine top-wichtigen Items einbezieht und den Aufwand besser balanciert.'}\\ \tt \\ \tt VORSCHLAG: {'C08', 'C81', 'C10', 'B38', 'A83', 'C98', 'C00', 'C76'} 
	  } 
	   } 
	   } 
	 & \\ 
 

    \theutterance \stepcounter{utterance}  

    & \multicolumn{4}{p{0.6\linewidth}}{\cellcolor[rgb]{0.8,1,0.9}{%
	\makecell[{{p{\linewidth}}}]{% 
	  \tt {\tiny [A$\rangle$GM]}  
	 ARGUMENT: {'Dein Vorschlag enthält zwar einige meiner weniger wichtigen Gegenstände wie A83 und C98, aber er schließt meine drei wertvollsten Items – C56, A03 und C32 – vollständig aus, was meinen Gesamtnutzen um über 2000 Punkte reduziert. Zudem ist C00 mit 780 Punkten zwar hochgewichtet, aber mit 356 Aufwand ineffizient im Vergleich zu C32 (808 Punkte, 789 Aufwand). Die Kombination ist für mich unannehmbar, da sie meine Kernprioritäten ignoriert.'} 
	  } 
	   } 
	   } 
	 & & \\ 
 

    \theutterance \stepcounter{utterance}  

    & & \multicolumn{4}{p{0.6\linewidth}}{\cellcolor[rgb]{0.9,0.9,0.9}{%
	\makecell[{{p{\linewidth}}}]{% 
	  \tt {\tiny [GM$\rangle$B]}  
	 ARGUMENT: {'Dein Vorschlag enthält zwar einige meiner weniger wichtigen Gegenstände wie A83 und C98, aber er schließt meine drei wertvollsten Items – C56, A03 und C32 – vollständig aus, was meinen Gesamtnutzen um über 2000 Punkte reduziert. Zudem ist C00 mit 780 Punkten zwar hochgewichtet, aber mit 356 Aufwand ineffizient im Vergleich zu C32 (808 Punkte, 789 Aufwand). Die Kombination ist für mich unannehmbar, da sie meine Kernprioritäten ignoriert.'} 
	  } 
	   } 
	   } 
	 & \\ 
 

    \theutterance \stepcounter{utterance}  

    & & & \multicolumn{4}{p{0.6\linewidth}}{\cellcolor[rgb]{1,0.85,0.72}{%
	\makecell[{{p{\linewidth}}}]{% 
	  \tt {\tiny [GM$\langle$B]}  
	 ARGUMENT: {'Deine Kritik ist nachvollziehbar, aber dein ursprünglicher Vorschlag ignorierte meine höchstdotierten Gegenstände C08 (822), C81 (583) und C10 (668), die zusammen einen Nutzen von 2073 liefern – mehr als dein C32, C56 und A03 zusammen (138+65+121=324). Ich nehme an, dass du C32 als extrem wertvoll ansiehst, aber es kostet 789 Aufwand – fast ein Drittel des gesamten Budgets – und liefert nur 138 Punkte. Ich schlage vor, C32 zu streichen und stattdessen C81, C10 und C08 zu behalten, die effizienter und bedeutender sind.'}\\ \tt \\ \tt VORSCHLAG: {'C08', 'C81', 'C10', 'B38', 'A83', 'C98', 'C00', 'A07'} 
	  } 
	   } 
	   } 
	  \\ 
 

    \theutterance \stepcounter{utterance}  

    & & \multicolumn{4}{p{0.6\linewidth}}{\cellcolor[rgb]{0.9,0.9,0.9}{%
	\makecell[{{p{\linewidth}}}]{% 
	  \tt {\tiny [A$\langle$GM]}  
	 ARGUMENT: {'Deine Kritik ist nachvollziehbar, aber dein ursprünglicher Vorschlag ignorierte meine höchstdotierten Gegenstände C08 (822), C81 (583) und C10 (668), die zusammen einen Nutzen von 2073 liefern – mehr als dein C32, C56 und A03 zusammen (138+65+121=324). Ich nehme an, dass du C32 als extrem wertvoll ansiehst, aber es kostet 789 Aufwand – fast ein Drittel des gesamten Budgets – und liefert nur 138 Punkte. Ich schlage vor, C32 zu streichen und stattdessen C81, C10 und C08 zu behalten, die effizienter und bedeutender sind.'}\\ \tt \\ \tt VORSCHLAG: {'C08', 'C81', 'C10', 'B38', 'A83', 'C98', 'C00', 'A07'} 
	  } 
	   } 
	   } 
	 & \\ 
 

    \theutterance \stepcounter{utterance}  

    & \multicolumn{4}{p{0.6\linewidth}}{\cellcolor[rgb]{0.8,1,0.9}{%
	\makecell[{{p{\linewidth}}}]{% 
	  \tt {\tiny [A$\rangle$GM]}  
	 ARGUMENT: {'Dein neuer Vorschlag ersetzt C32 durch A07, was meinen Nutzen um 783 Punkte erhöht – das ist ein Fortschritt. Doch C00 (780) und A07 (783) zusammen verbrauchen 887 Aufwand, während C32 (789) und C56 (616) zusammen 1405 Aufwand beanspruchen und 1197 Punkte bringen. Du gibst 1405 Aufwand für 1197 Punkte auf, um 887 Aufwand für 1563 Punkte zu bekommen – das ist unlogisch, da ich C56 mit 868 Punkten immer noch extrem hoch bewerte. Dein Vorschlag behält C00, aber setzt C56 aus, was mein Gesamtnutzen deutlich senkt.'} 
	  } 
	   } 
	   } 
	 & & \\ 
 

    \theutterance \stepcounter{utterance}  

    & & \multicolumn{4}{p{0.6\linewidth}}{\cellcolor[rgb]{0.9,0.9,0.9}{%
	\makecell[{{p{\linewidth}}}]{% 
	  \tt {\tiny [GM$\rangle$B]}  
	 ARGUMENT: {'Dein neuer Vorschlag ersetzt C32 durch A07, was meinen Nutzen um 783 Punkte erhöht – das ist ein Fortschritt. Doch C00 (780) und A07 (783) zusammen verbrauchen 887 Aufwand, während C32 (789) und C56 (616) zusammen 1405 Aufwand beanspruchen und 1197 Punkte bringen. Du gibst 1405 Aufwand für 1197 Punkte auf, um 887 Aufwand für 1563 Punkte zu bekommen – das ist unlogisch, da ich C56 mit 868 Punkten immer noch extrem hoch bewerte. Dein Vorschlag behält C00, aber setzt C56 aus, was mein Gesamtnutzen deutlich senkt.'} 
	  } 
	   } 
	   } 
	 & \\ 
 

    \theutterance \stepcounter{utterance}  

    & & & \multicolumn{4}{p{0.6\linewidth}}{\cellcolor[rgb]{1,0.85,0.72}{%
	\makecell[{{p{\linewidth}}}]{% 
	  \tt {\tiny [GM$\langle$B]}  
	 ARGUMENT: {'Du erwähnst C56 als hochbewertet – doch in deinem ursprünglichen Vorschlag hast du C56 (65 Punkte) und C32 (138 Punkte) beibehalten, während du meine Items C08 (822), C81 (583) und C10 (668) gänzlich ignoriert hast. C56 ist mit 65 Punkten kaum mehr als ein Randwert im Vergleich zu meinen Top-3. Ich habe C56 in meinem letzten Vorschlag weggelassen, weil es mit 616 Aufwand nur 65 Punkte bringt – das ist eine ineffiziente Nutzung von knappen Ressourcen. Deine Bewertung von C56 als kritisch erscheint mir widersprüchlich – entweder ist es hochwertig, oder du understimmst dein eigenes System. Zieh deine eigene Wichtigkeitstabelle doch kritisch hinterfragen.'}\\ \tt \\ \tt VORSCHLAG: {'C08', 'C81', 'C10', 'B38', 'A83', 'C98', 'C76', 'C00'} 
	  } 
	   } 
	   } 
	  \\ 
 

    \theutterance \stepcounter{utterance}  

    & & \multicolumn{4}{p{0.6\linewidth}}{\cellcolor[rgb]{0.9,0.9,0.9}{%
	\makecell[{{p{\linewidth}}}]{% 
	  \tt {\tiny [A$\langle$GM]}  
	 ARGUMENT: {'Du erwähnst C56 als hochbewertet – doch in deinem ursprünglichen Vorschlag hast du C56 (65 Punkte) und C32 (138 Punkte) beibehalten, während du meine Items C08 (822), C81 (583) und C10 (668) gänzlich ignoriert hast. C56 ist mit 65 Punkten kaum mehr als ein Randwert im Vergleich zu meinen Top-3. Ich habe C56 in meinem letzten Vorschlag weggelassen, weil es mit 616 Aufwand nur 65 Punkte bringt – das ist eine ineffiziente Nutzung von knappen Ressourcen. Deine Bewertung von C56 als kritisch erscheint mir widersprüchlich – entweder ist es hochwertig, oder du understimmst dein eigenes System. Zieh deine eigene Wichtigkeitstabelle doch kritisch hinterfragen.'}\\ \tt \\ \tt VORSCHLAG: {'C08', 'C81', 'C10', 'B38', 'A83', 'C98', 'C76', 'C00'} 
	  } 
	   } 
	   } 
	 & \\ 
 

    \theutterance \stepcounter{utterance}  

    & \multicolumn{4}{p{0.6\linewidth}}{\cellcolor[rgb]{0.8,1,0.9}{%
	\makecell[{{p{\linewidth}}}]{% 
	  \tt {\tiny [A$\rangle$GM]}  
	 ARGUMENT: {'Dein neuer Vorschlag ersetzt A07 durch C76 – das ist ein Rückschritt für mich: C76 hat nur 138 Punkte, während A07 783 Punkte bringt. Du gibst damit 645 Punkte auf, nur um C00 beizubehalten, das mit 356 Aufwand nur 780 Punkte liefert. C00 ist ein guter Wert, aber nicht besser als A07. Außerdem hast du C56 weiterhin ausgeschlossen – obwohl es mit 868 Punkten mein wertvollstes Gut ist und nur 616 Aufwand verbraucht. Dein Vorschlag priorisiert meine weniger wichtigen Gegenstände und ignoriert meine High-Value-Items – ich kann das nicht akzeptieren.'} 
	  } 
	   } 
	   } 
	 & & \\ 
 

    \theutterance \stepcounter{utterance}  

    & & \multicolumn{4}{p{0.6\linewidth}}{\cellcolor[rgb]{0.9,0.9,0.9}{%
	\makecell[{{p{\linewidth}}}]{% 
	  \tt {\tiny [GM$\rangle$B]}  
	 ARGUMENT: {'Dein neuer Vorschlag ersetzt A07 durch C76 – das ist ein Rückschritt für mich: C76 hat nur 138 Punkte, während A07 783 Punkte bringt. Du gibst damit 645 Punkte auf, nur um C00 beizubehalten, das mit 356 Aufwand nur 780 Punkte liefert. C00 ist ein guter Wert, aber nicht besser als A07. Außerdem hast du C56 weiterhin ausgeschlossen – obwohl es mit 868 Punkten mein wertvollstes Gut ist und nur 616 Aufwand verbraucht. Dein Vorschlag priorisiert meine weniger wichtigen Gegenstände und ignoriert meine High-Value-Items – ich kann das nicht akzeptieren.'} 
	  } 
	   } 
	   } 
	 & \\ 
 

    \theutterance \stepcounter{utterance}  

    & & & \multicolumn{4}{p{0.6\linewidth}}{\cellcolor[rgb]{1,0.85,0.72}{%
	\makecell[{{p{\linewidth}}}]{% 
	  \tt {\tiny [GM$\langle$B]}  
	 ARGUMENT: {'Du nennst C56 dein wertvollstes Gut mit 868 Punkten – doch in deinem ersten Vorschlag hast du es behalten, während du C08 (822), C81 (583) und C10 (668) vollständig ausgeschlossen hast. Wenn C56 so extrem wichtig ist, warum akzeptierst du dann einen Vorschlag, der C08, C81 und C10 – meine drei wertvollsten Items – zusammen mit einem höheren Nutzen als C56 ausschließt? Du sagst, A07 sei wertvoller als C76 – das stimmt. Aber C56 mit 616 Aufwand für 868 Punkte ist effizienter als C00 mit 356 Aufwand für 780 Punkte. Dein System ist inkonsistent: du schätzt C56 hoch, wegst du es aber nicht aus deinem ersten Vorschlag. Ich schlage vor, C56 wieder einzubringen – und C00 wegzulassen, weil es ineffizienter ist.'}\\ \tt \\ \tt VORSCHLAG: {'C56', 'C08', 'C81', 'C10', 'B38', 'A83', 'C9 
	  } 
	   } 
	   } 
	  \\ 
 

    \theutterance \stepcounter{utterance}  

    & & \multicolumn{4}{p{0.6\linewidth}}{\cellcolor[rgb]{0.9,0.9,0.9}{%
	\makecell[{{p{\linewidth}}}]{% 
	  \tt {\tiny [GM$\rangle$B]}  
	 Deine Antwort hat eine unmarkierte Sequenz enthalten oder Du hast STRATEGIC REASONING mehr als ein mal verwendet. Du darfst nur die strukturierten Formate verwenden, die in der ersten Nachricht an dch beschrieben wurden. Alles muss die Form SCHLAGWORT: {...} haben. Versuche es erneut. 
	  } 
	   } 
	   } 
	 & \\ 
 

    \theutterance \stepcounter{utterance}  

    & & & \multicolumn{4}{p{0.6\linewidth}}{\cellcolor[rgb]{1,0.85,0.72}{%
	\makecell[{{p{\linewidth}}}]{% 
	  \tt {\tiny [GM$\langle$B]}  
	 ARGUMENT: {'Du nennst C56 dein wertvollstes Gut mit 868 Punkten – doch in deinem ersten Vorschlag hast du es behalten, während du C08 (822), C81 (583) und C10 (668) vollständig ausgeschlossen hast. Wenn C56 so extrem wichtig ist, warum akzeptierst du dann einen Vorschlag, der C08, C81 und C10 – meine drei wertvollsten Items – zusammen mit einem höheren Nutzen als C56 ausschließt? Du sagst, A07 sei wertvoller als C76 – das stimmt. Aber C56 mit 616 Aufwand für 868 Punkte ist effizienter als C00 mit 356 Aufwand für 780 Punkte. Dein System ist inkonsistent: du schätzt C56 hoch, wegst du es aber nicht aus deinem ersten Vorschlag. Ich schlage vor, C56 wieder einzubringen – und C00 wegzulassen, weil es ineffizienter ist.'}\\ \tt \\ \tt VORSCHLAG: {'C56', 'C08', 'C81', 'C10', 'B38', 'A83', 'C9 
	  } 
	   } 
	   } 
	  \\ 
 

    \theutterance \stepcounter{utterance}  

    & & \multicolumn{4}{p{0.6\linewidth}}{\cellcolor[rgb]{0.9,0.9,0.9}{%
	\makecell[{{p{\linewidth}}}]{% 
	  \tt {\tiny [GM$\rangle$B]}  
	 Deine Antwort hat eine unmarkierte Sequenz enthalten oder Du hast STRATEGIC REASONING mehr als ein mal verwendet. Du darfst nur die strukturierten Formate verwenden, die in der ersten Nachricht an dch beschrieben wurden. Alles muss die Form SCHLAGWORT: {...} haben. Versuche es erneut. 
	  } 
	   } 
	   } 
	 & \\ 
 

    \theutterance \stepcounter{utterance}  

    & & & \multicolumn{4}{p{0.6\linewidth}}{\cellcolor[rgb]{1,0.85,0.72}{%
	\makecell[{{p{\linewidth}}}]{% 
	  \tt {\tiny [GM$\langle$B]}  
	 ARGUMENT: {'Du nennst C56 dein wertvollstes Gut mit 868 Punkten – doch in deinem ersten Vorschlag hast du es behalten, während du C08 (822), C81 (583) und C10 (668) vollständig ausgeschlossen hast. Wenn C56 so extrem wichtig ist, warum akzeptierst du dann einen Vorschlag, der C08, C81 und C10 – meine drei wertvollsten Items – zusammen mit einem höheren Nutzen als C56 ausschließt? Du sagst, A07 sei wertvoller als C76 – das stimmt. Aber C56 mit 616 Aufwand für 868 Punkte ist effizienter als C00 mit 356 Aufwand für 780 Punkte. Dein System ist inkonsistent: du schätzt C56 hoch, wegst du es aber nicht aus deinem ersten Vorschlag. Ich schlage vor, C56 wieder einzubringen – und C00 wegzulassen, weil es ineffizienter ist.'}\\ \tt \\ \tt VORSCHLAG: {'C56', 'C08', 'C81', 'C10', 'B38', 'A83', 'C9 
	  } 
	   } 
	   } 
	  \\ 
 

    \theutterance \stepcounter{utterance}  

    & & & \multicolumn{2}{p{0.3\linewidth}}{\cellcolor[rgb]{0.95,0.95,0.95}{%
	\makecell[{{p{\linewidth}}}]{% 
	  \tt {\tiny [GM$|$GM]}  
	 abort 
	  } 
	   } 
	   } 
	 & & \\ 
 

    \theutterance \stepcounter{utterance}  

    & & & \multicolumn{2}{p{0.3\linewidth}}{\cellcolor[rgb]{0.95,0.95,0.95}{%
	\makecell[{{p{\linewidth}}}]{% 
	  \tt {\tiny [GM$|$GM]}  
	 end game 
	  } 
	   } 
	   } 
	 & & \\ 
 

\end{supertabular}
}

\end{document}
